% 02-ky-hieu.tex - The book's notation system
%
% Scope: vectors, matrices, dimensions, the recurrence equation, and why
% one notation is needed for six conflicting papers.

\section{Hệ ký hiệu}
\label{sec:ch01-ky-hieu}

Trước khi viết phương trình đầu tiên, tôi cần chốt một bộ ký hiệu dùng cho
toàn bộ cuốn sách. Đây không phải là việc trang trí. Sáu bài báo gốc mà cuốn
sách này dựa vào dùng sáu bộ ký hiệu khác nhau và không bộ nào tương thích
với bộ nào. LSTM~1997 gọi ma trận trọng số hồi quy là $\mat{U}$; Pascanu~2013
gọi nó là $\mat{W}_{\text{rec}}$; Transformer~2017 không có ma trận hồi quy
nào cả nhưng có ba ma trận attention đều tên là $W$. Nếu mỗi chương dùng ký
hiệu của bài báo mà nó đang đọc, bạn sẽ phải học lại bảng chữ cái ở mỗi
chương, và tôi sẽ không thể đặt hai phương trình cạnh nhau để so sánh; đó
lại chính là việc cuốn sách này tồn tại để làm.

Tôi chọn một bộ ký hiệu theo ba nguyên tắc. Một: nhất quán toàn sách; mỗi
đại lượng có đúng một tên, mỗi tên có đúng một nghĩa. Hai: gần với quy ước
hiện đại; bạn ra ngoài đọc paper năm 2026 sẽ gặp $\vect{h}_t$ cho trạng thái
ẩn, không gặp $\vect{s}_t$ như paper RNN đầu tiên. Ba: tự mô tả; $\mat{W}_{xh}$
tự nói ``từ $x$ tới $h$'', không cần bạn tra bảng.

\subsection{Đại lượng}

Vector được in đậm chữ thường, ma trận in đậm chữ hoa. Tại mỗi
\tn{bước thời gian}{timestep} $t = 1, 2, \dots, T$, mạng nhận một vector đầu
vào $\vect{x}_t \in \R^{d_x}$, giữ một trạng thái ẩn $\vect{h}_t \in \R^{d_h}$,
và sinh một vector đầu ra $\vect{y}_t \in \R^{d_y}$. Ký hiệu $d_x, d_h, d_y$
là số chiều của từng không gian; chúng là hằng số với một mô hình cho trước
và không phụ thuộc vào $T$. Trạng thái ẩn ban đầu $\vect{h}_0$ luôn là vector
không: mạng bắt đầu với một tờ giấy trắng.

Đầu vào $\vect{x}_t$ có thể là một từ được nhúng trong không gian
embedding, một vector đặc trưng được trích từ khung hình thứ $t$ của video,
hay một điểm dữ liệu trong chuỗi thời gian; RNN không quan tâm. Điều nó
làm với đầu vào đó là như nhau ở mọi bài toán, và đó là thứ ta quan tâm
trong chương này.

\subsection{Phương trình hồi quy}

Một RNN được định nghĩa bởi hai phương trình:

\begin{align}
  \vect{h}_t &= \tanh\!\bigl(\mat{W}_{xh} \vect{x}_t
               + \mat{W}_{hh} \vect{h}_{t-1}
               + \vect{b}_h\bigr), \label{eq:rnn-hidden} \\
  \vect{y}_t &= \mat{W}_{hy} \vect{h}_t + \vect{b}_y. \label{eq:rnn-output}
\end{align}

Phương trình~\eqref{eq:rnn-hidden} là trái tim của RNN. Tại mỗi bước $t$, nó
lấy đầu vào hiện tại $\vect{x}_t$ và trạng thái ẩn trước đó $\vect{h}_{t-1}$,
nhân với hai \tn{ma trận trọng số}{weight matrix} $\mat{W}_{xh}$ (input
$\to$ hidden) và $\mat{W}_{hh}$ (hidden $\to$ hidden), cộng
\tn{bias}{bias} $\vect{b}_h$, rồi đẩy qua hàm $\tanh$. Kết quả là trạng thái
ẩn mới $\vect{h}_t$; một vector tổng kết mọi thứ mạng đã thấy từ bước~1 đến
bước~$t$.

Phương trình~\eqref{eq:rnn-output} đọc trạng thái ẩn ra thành đầu ra
$\vect{y}_t$ qua một phép biến đổi tuyến tính nữa. Trong bài toán phân loại,
$\vect{y}_t$ sẽ được đưa qua softmax; trong bài toán hồi quy, nó là đầu ra
trực tiếp. Chương này dùng hồi quy vì nó giữ cho đạo hàm đơn giản; mỗi bước
thêm một lớp softmax là thêm một tầng Jacobian không cần thiết cho mục tiêu
hiểu BPTT.

\subsection{Tháo cuộn theo thời gian}

Một RNN có vẻ tuần hoàn: $\vect{h}_t$ xuất hiện ở cả vế trái lẫn vế phải của
\eqref{eq:rnn-hidden}. Nhưng với một chuỗi có độ dài $T$ cố định, ta có thể
\textbf{tháo cuộn} nó thành một mạng truyền thẳng sâu $T$ tầng, mỗi tầng dùng
chung một bộ trọng số. Hình dung: tại $t=1$, mạng tính $\vect{h}_1$ từ
$\vect{x}_1$ và $\vect{h}_0$; tại $t=2$, nó tính $\vect{h}_2$ từ $\vect{x}_2$
và $\vect{h}_1$; và cứ thế đến $T$.

Chính hình ảnh tháo cuộn này là thứ làm cho BPTT khả thi. Một khi chuỗi đã
được trải phẳng thành một đồ thị tính toán không có chu trình, ta có thể áp
dụng đúng thuật toán lan truyền ngược thông thường; chỉ khác là gradient
phải đi qua cùng một tập trọng số $\mat{W}_{hh}$ nhiều lần, mỗi lần tại một
bước thời gian. Chính sự lặp lại đó là nguồn gốc của mọi hiện tượng mà phần~II
của cuốn sách sẽ mổ xẻ.

\subsection{Ký hiệu gradient}

Cuốn sách này làm việc với gradient ở dạng \textbf{ma trận Jacobi}. Nếu
$\vect{y} \in \R^m$ và $\vect{x} \in \R^n$, thì Jacobian
$\jac{\vect{y}}{\vect{x}}$ là một ma trận $m \times n$ với phần tử
$(i,j)$ là $\partial y_i / \partial x_j$. Với một \tn{hàm mất mát}{loss
function} vô hướng $L$, gradient của $L$ theo một ma trận $\mat{W}$ là một
ma trận cùng kích thước:

\[
  \pd{L}{\mat{W}} = \begin{bmatrix}
    \frac{\partial L}{\partial W_{11}} & \cdots & \frac{\partial L}{\partial W_{1n}} \\
    \vdots & \ddots & \vdots \\
    \frac{\partial L}{\partial W_{m1}} & \cdots & \frac{\partial L}{\partial W_{mn}}
  \end{bmatrix}.
\]

Dạng Jacobian có một lợi thế lớn: quy tắc dây chuyền trở thành phép nhân ma
trận. Nếu $\vect{z} = f(\vect{y})$ và $\vect{y} = g(\vect{x})$, thì

\[
  \jac{\vect{z}}{\vect{x}} = \jac{\vect{z}}{\vect{y}} \; \jac{\vect{y}}{\vect{x}}.
\]

Không cần nhớ khi nào thì nhân với chuyển vị, không cần tranh luận
\emph{denominator layout} hay \emph{numerator layout}; Jacobian tự xử lý
chiều. Từ chương này đến hết chương~4, mọi dẫn xuất sẽ được viết ở dạng này.
Khi đến chương~7 với Transformer, ta sẽ chuyển sang dạng vector vì bản thân
bài báo Vaswani làm vậy; phụ lục~A ghi lại ánh xạ giữa hai dạng cho từng bài
báo.

\subsection{Ký hiệu cho BPTT}

Để dẫn xuất BPTT, tôi tách pre-activation khỏi hidden state. Gọi
$\vect{a}_t$ là vector \textbf{trước} khi qua $\tanh$:

\begin{equation}
  \vect{a}_t = \mat{W}_{xh} \vect{x}_t
              + \mat{W}_{hh} \vect{h}_{t-1}
              + \vect{b}_h,
  \qquad
  \vect{h}_t = \tanh(\vect{a}_t).
  \label{eq:rnn-preact}
\end{equation}

Ký hiệu $\vect{a}_t$ không có trong các bài báo gốc; họ viết thẳng
$\vect{h}_t = \tanh(\mat{W}_{xh} \vect{x}_t + \mat{W}_{hh} \vect{h}_{t-1}
+ \vect{b}_h)$. Nhưng khi dẫn xuất BPTT, tách $\vect{a}_t$ ra giúp cô lập
đạo hàm của $\tanh$ thành một bước riêng, và bước đó là ma trận đường chéo
$\diag(1 - \vect{h}_t^2)$. Sự tách biệt này sẽ trả giá trong mục~4 khi ta
tính gradient qua từng cổng.
